\item 
\textbf{HydroSense Tech }    \hfill     
\hfill  \textit{Sep 2017 -- Present (Part Time, remote); 
\item  Co-founder / Senior Data Scientist (Sensor Data)\hfill May 2025 -- Sep 2025 (Full Time, Singapore)}
\vspace*{-3pt}
\begin{itemize}[leftmargin=0.2in, topsep=1pt]
  \item Co-founded an IoT sensing startup and owned end-to-end sensor data curation and analysis, setting data-quality KPIs and dashboards to drive weekly DC/DA workflows and improve field-device reliability.
  \item Built robust sensor data pipelines from raw acquisition to cleaned, analyzable datasets, including time sync, calibration, filtering, and feature extraction in Python/MATLAB, enabling scalable analyses across multi-site deployments.
  \item Developed visualization and root-cause analysis tooling (range/time plots, anomaly triage, drift tracking) to debug edge cases and biases, accelerating iteration with DSP/ML engineers and improving model-ready data quality.
  \item Implemented automation for repetitive labeling and validation tasks, including heuristics-based pre-labeling and sanity checks, reducing manual review effort while maintaining strict integrity and traceability of curated datasets.
  \item \textbf{Patents:} Co-inventor on 3 granted patents covering rainfall sensing hardware, sensor calibration, and ML-based drift compensation / EMI mitigation algorithms for industrial monitoring applications.
\end{itemize}

\vspace*{-1pt}
\item \textbf{Seno Medical}\\
\textit{Senior Data Scientist (Sensor Data)} \hfill Buffalo, NY \quad \textit{May 2024 -- Aug 2024}
\vspace*{-3pt}
\begin{itemize}[leftmargin=0.2in, topsep=1pt]
  \item Owned data curation and exploratory analysis for multimodal sensor datasets, defining quality gates, labeling guidelines, and weekly review cadences to surface gaps, biases, and edge cases impacting downstream models.
  \item Built publication-quality visualizations and diagnostic reports for stakeholders, translating complex signal behaviors into actionable insights for target detection, tracking, and classification workflows across DSP and ML teams.
  \item Designed reproducible processing pipelines in Python (NumPy, SciPy, scikit-learn) with version control and documentation, enabling consistent cleaning, transformation, and feature engineering across multiple data sources.
  \item Partnered with engineers to validate end-to-end RF/ML pipelines and run RCA on failures, using structured experiments and statistical checks to prioritize fixes and improve system-level performance.
\end{itemize}

\vspace*{-1pt}
\item \textbf{Roswell Park Comprehensive Cancer Center}\\
\textit{Research Scientist} \hfill Buffalo, NY \quad \textit{May 2023 -- Aug 2023}
\vspace*{-3pt}
\begin{itemize}[leftmargin=0.2in, topsep=1pt]
  \item Designed sensor/imaging data pipelines aligning raw acquisition with structured metadata, including de-identification, schema normalization, and automated QC, producing clean datasets for analysis and ML model development.
  \item Built scalable dataset generation jobs (filtering, sampling, augmentation, tagging) and error-analysis scripts to improve coverage of rare edge cases, enabling robust training and evaluation under distribution shifts.
  \item Implemented signal processing and ML workflows (denoising, calibration, representation learning) in PyTorch and Python, improving SNR and stability for downstream computer-vision and quantitative biomarker tasks.
  \item Automated experiments and reproducibility with standardized configs, logging, and comparison notebooks, enabling faster iteration and clearer RCA across model variants, preprocessing choices, and data-quality interventions.
\end{itemize}

